% ID: SEL-D0063-OB-E04
% ESTADO: cotejada. Segunda pasada visual y borrador compilado revisados.
% Fuente: Selectividad/sel_2010_matematicas_aplicadas/mat_apli/mat_apli_junio/mat_apli_junio_rb/EXAMEN-ANDALUCÍA-6.pdf
% SHA-256: 0de2f7d077a3b5d5c0dbca886cdb363082356411843369c3bce54ff7cda08abc
% Página: 2; opción B; ejercicio original 4.

El peso de los sacos de patatas de una cooperativa es una variable aleatoria Normal con desviación típica $0.25$ kg. El agente de ventas de esa cooperativa afirma que el peso medio de los sacos no baja de 5 kg.

Se desea contrastar estadísticamente esta hipótesis. Para ello se toma una muestra aleatoria de 20 sacos y se obtiene que su peso medio es de $4.8$ kg.
\begin{enumerate}[label=\alph*)]
\item \textbf{(0.5 puntos)} Determine las hipótesis del contraste que se plantea en este enunciado.
\item \textbf{(1 punto)} Halle la región crítica de este contraste para $\alpha=0.01$.
\item \textbf{(1 punto)} Con los datos de la muestra tomada, \textquestiondown{}puede decirse que existe evidencia estadística suficiente para rechazar la hipótesis del agente de ventas de la cooperativa, al nivel de significación $\alpha=0.01$?
\end{enumerate}
